\documentclass[a4paper, serif, 10pt]{moderncv}

%% moderncv
% color and style
\moderncvcolor{black}
\moderncvstyle{banking}

% renew moderncv command to adapt for CJK colon
\renewcommand*{\cvitem}[3][.25em]{%
  \ifstrempty{#2}{}{\hintstyle{#2}：}{#3}%
  \par\addvspace{#1}}

\renewcommand*{\cvitemwithcomment}[4][.25em]{%
  \savebox{\cvitemwithcommentmainbox}{\ifstrempty{#2}{}{\hintstyle{#2}：}#3}%
  \setlength{\cvitemwithcommentmainlength}{\widthof{\usebox{\cvitemwithcommentmainbox}}}%
  \setlength{\cvitemwithcommentcommentlength}{\maincolumnwidth-\separatorcolumnwidth-\cvitemwithcommentmainlength}%
  \begin{minipage}[t]{\cvitemwithcommentmainlength}\usebox{\cvitemwithcommentmainbox}\end{minipage}%
  \hfill% fill of \separatorcolumnwidth
  \begin{minipage}[t]{\cvitemwithcommentcommentlength}\raggedleft\small\itshape#4\end{minipage}%
  \par\addvspace{#1}}

%% page layout/margins
\usepackage[top=2.5cm, bottom=2.5cm, left=1.5cm, right=1.5cm]{geometry}

\nopagenumbers{}

%% Babel config for English language
\usepackage[english]{babel}

%% fontspec
\usepackage{fontspec}

\IfFontExistsTF{Linux Libertine}{
  \setmainfont[Ligatures={TeX, Common}, Numbers=Lining, ItalicFont=Linux Libertine]{Linux Libertine}
}{}
\IfFontExistsTF{Linux Libertine O}{
  \setmainfont[Ligatures={TeX, Common}, Numbers=Lining, ItalicFont=Linux Libertine O]{Linux Libertine O}
}{}

%% CTeX
% CJK support, used to show CJK characters in the resume
%
% - heading=false: disable ctex heading
% - punct=kaiming: use kaiming punctuations styles for CJK
% - scheme=plain: use plain scheme, do not override \normalsize font size
% - space=auto: space settings for CJK characters
%
% CJK fonts are resolved per locale language, e.g. Noto Serif/Sans CJK JP
% for Japanese resumes, with fallback to the next region if not installed
%
% ref:
% - http://ctan.mirrorcatalogs.com/language/chinese/ctex/ctex.pdf
\usepackage[UTF8, heading=false, punct=kaiming, scheme=plain, space=auto]{ctex}

\IfFontExistsTF{Noto Serif CJK HK}{
  \setCJKmainfont{Noto Serif CJK HK}
}{
\IfFontExistsTF{Noto Serif CJK TC}{
  \setCJKmainfont{Noto Serif CJK TC}
}{}}
\IfFontExistsTF{Noto Sans CJK HK}{
  \setCJKsansfont{Noto Sans CJK HK}
}{
\IfFontExistsTF{Noto Sans CJK TC}{
  \setCJKsansfont{Noto Sans CJK TC}
}{}}

%% line spacing
\usepackage{setspace}
\setstretch{1.125}

%% URL styling - use normal text instead of monospace
\urlstyle{same}

%% Patch \httplink and \httpslink to handle full URLs
% The original moderncv macros always prepend http:// or https:// to URLs.
% This causes issues when the URL already has a protocol (e.g., https://example.com)
% resulting in malformed URLs like https://https://example.com.
% This patch checks if the URL already contains "://" and uses it directly if so.
% Note: We use \str_set:Nx (x-expansion) to expand macros like \@homepage first.
\ExplSyntaxOn
\str_new:N \g_yamlresume_url_str

\RenewDocumentCommand{\httpslink}{O{}m}{
  \str_set:Nx \g_yamlresume_url_str {#2}
  \str_if_in:NnTF \g_yamlresume_url_str {://}
    { \url{#2} }
    { \url{https://#2} }
}

\RenewDocumentCommand{\httplink}{O{}m}{
  \str_set:Nx \g_yamlresume_url_str {#2}
  \str_if_in:NnTF \g_yamlresume_url_str {://}
    { \url{#2} }
    { \url{http://#2} }
}
\ExplSyntaxOff

\name{陳家豪}{}
\title{高級財務分析師 | CFA 特許金融分析師}
\phone[mobile]{+852 9123 4567}
\email{ka.ho.chan@example.com}
\homepage{https://www.chan-kaho-finance.com}

\address{中國香港，香港島，香港，香港中環皇后大道中99號}{}{}

\extrainfo{{\small \faLinkedin}\ \href{https://www.linkedin.com/in/chan-ka-ho-finance}{@chan-ka-ho-finance} {} {} {} • {} {} {} 
{\small \faMedium}\ \href{https://medium.com/@chan.kaho.finance}{@chan.kaho.finance}}

\begin{document}

\maketitle

\section{簡介}

\cvline{}{具備 7 年投資銀行與資產管理經驗的高級財務分析師，專注於財務建模、估值分析、預算預測及管理層決策支援。
%
\begin{itemize}
\item 精通三表模型、DCF、可比公司分析及情境分析，能將複雜數據轉化為可執行洞察
\item 熟悉香港上市公司財務披露、IFRS 及 ESG 報告要求
\item 擅長以 Excel、SQL、Python 及 Power BI 建立自動化儀表板，提升月度結算與差異分析效率
\item 持有 CFA 特許資格，具備優秀的中英文溝通及跨部門協作能力
\end{itemize}}
\section{教育背景}

\cventry{2014年9月–2018年6月}
        {學士，計量金融學，成績：3.7/4.0}
        {香港中文大學}
        {\url{https://www.cuhk.edu.hk}}
        {}
        {\begin{itemize}
\item 主修計量金融，成績排名前 10\%，連續兩年獲院長榮譽錄
\item 完成企業估值、投資組合管理及固定收益證券等專題研究
\item 擔任金融學會學術幹事，組織行業分享會及模擬投資比賽
\end{itemize}
\textbf{課程}：財務管理、投資分析與資產定價、衍生工具與風險管理、計量方法與統計推論、會計學原理、企業估值、金融市場與機構、程式設計與數據分析}
\section{工作經歷}

\cventry{2021年3月至今}
        {高級財務分析師}
        {香港交易所}
        {\url{https://www.hkex.com.hk}}
        {}
        {\begin{itemize}
\item 負責集團月度及季度財務分析，涵蓋收入、成本、資本開支及現金流，支援管理層制定年度預算
\item 建立並維護三表財務模型、DCF 及可比公司估值模型，用於新業務投資及策略評估
\item 主導差異分析與滾動預測，將預測週期由 4 週縮短至 10 個工作天
\item 與業務、風險及合規部門協作，編製董事會及監管機構匯報材料
\item 指導初級分析師，制定財務建模及數據視覺化最佳實務
\end{itemize}
\textbf{關鍵字}：財務建模、估值分析、預算預測、差異分析、管理層匯報、團隊指導}

\cventry{2018年7月–2021年2月}
        {財務分析師}
        {恒生銀行}
        {\url{https://www.hangseng.com}}
        {}
        {\begin{itemize}
\item 支援零售銀行業務的預算編製、月度結算及績效追蹤，管理年度預算逾 8 億港元
\item 分析產品盈利能力及客戶分群表現，提出定價及資源配置建議
\item 開發 Excel VBA 及 Power BI 自動化報表，減少每月手動處理時間約 60\%
\item 參與 IFRS 9 撥備模型校驗及監管報表準備工作
\end{itemize}
\textbf{關鍵字}：預算管理、績效分析、產品盈利、監管報表、Power BI、Excel VBA}
\section{語言}

\cvline{粵語}{母語或雙語水平 \hfill \textbf{關鍵字}：母語}
\cvline{英語}{完全專業水平 \hfill \textbf{關鍵字}：IELTS 7.5、商務簡報}
\cvline{普通話}{完全專業水平 \hfill \textbf{關鍵字}：普通話水平測試二級甲等、跨區業務溝通}
\section{專業技能}

\cvline{財務分析與建模}{專家 \hfill \textbf{關鍵字}：三表財務模型、DCF 估值、可比公司分析、情境與敏感度分析}
\cvline{數據分析與視覺化}{高級 \hfill \textbf{關鍵字}：Excel VBA、SQL、Python、Power BI、Tableau}
\cvline{會計與報告}{高級 \hfill \textbf{關鍵字}：IFRS、管理層報告、預算與預測、差異分析}
\cvline{金融工具}{高級 \hfill \textbf{關鍵字}：Bloomberg Terminal、Capital IQ、FactSet、SAP ERP}
\section{榮譽}

\cventry{2022年12月}
        {香港交易所}
        {最佳財務分析項目獎}
        {}
        {}
        {表揚在集團新業務投資評估中建立自動化估值模型，將分析週期縮短 35\%，並提升董事會決策資料的準確性。}

\cventry{2017年5月}
        {香港中文大學}
        {院長榮譽錄}
        {}
        {}
        {因學術成績優異而獲頒，連續兩年位列商學院前 10\%。}
\section{證書}

\cventry{2022年8月}
        {CFA 協會}
        {CFA 特許金融分析師資格}
        {\url{https://www.cfainstitute.org}}
        {}
        {}

\cventry{2020年11月}
        {全球風險管理專業人員協會}
        {金融風險管理師（FRM）}
        {\url{https://www.garp.org}}
        {}
        {}
\section{出版刊物}

\cventry{2023年6月}
        {香港上市公司 ESG 披露與估值相關性研究}
        {香港財務分析師學會期刊}
        {\url{https://www.hksfa.org/publications/esg-valuation}}
        {}
        {\begin{itemize}
\item 分析 2018 至 2022 年香港上市公司 ESG 評分與市盈率、資金成本的關係
\item 發現高 ESG 評分企業在市場波動期間展現較低資金成本
\item 提出適用於本地市場的 ESG 整合估值框架
\end{itemize}}
\section{項目}

\cventry{2023年1月–2023年6月}
        {結合 ESG 評分、財務預測及情境分析的可配置估值模型，供投資研究團隊使用。}
        {港股 ESG 估值整合模型}
        {\url{https://github.com/chan-ka-ho-finance/esg-valuation-model}}
        {}
        {\begin{itemize}
\item 以 Python 串接公開 ESG 數據及財務報表，自動計算調整後資金成本
\item 建立行業中性評分及敏感度分析，支援投資委員會討論
\item 模型獲研究部門採用於 20 間恒生指數成分股之初步篩選
\end{itemize}
\textbf{關鍵字}：ESG、估值模型、Python、情境分析}

\cventry{2020年2月–2020年11月}
        {為零售銀行業務設計的預算監控與差異分析儀表板，整合多個資料來源並自動更新。}
        {零售銀行預算自動化儀表板}
        {\url{https://github.com/chan-ka-ho-finance/budget-dashboard}}
        {}
        {\begin{itemize}
\item 使用 SQL 及 Power BI 建立產品、渠道及客戶分群維度的預算追蹤
\item 將每月管理報表編製時間由 3 天縮短至半天
\item 增設異常警示及下鑽分析，協助業務主管快速定位偏差原因
\end{itemize}
\textbf{關鍵字}：Power BI、SQL、預算管理、差異分析}
\section{興趣愛好}

\cvline{金融科技}{區塊鏈、量化交易、數據視覺化}
\cvline{運動}{跑步、羽毛球、遠足}
\section{志願服務}

\cventry{2019年9月–2023年12月}
        {社區理財教育義工}
        {香港財務策劃師學會}
        {\url{https://www.ifphk.org}}
        {}
        {\begin{itemize}
\item 為中學生及基層家庭主持理財教育工作坊，主題包括預算管理、儲蓄及風險保障
\item 協助設計互動教材及案例，累計服務超過 300 名參加者
\item 與社區中心合作統籌年度理財同樂日，協調義工及活動流程
\end{itemize}}
    
\end{document}